\scanpage{18}
\mspage{\handwritten{\deleted{\unsure{48|18}}4}}
\handwritten{l'}exposé \deleted{ultérieur} \inserted{\handwritten{$\overline{\underline{\mathrm{XIV}}}$}} \deleted{que la réciproque est vraie} \inserted{\handwritten{que la réciproque est vraie}} lorsque $S$ est \deleted{\unsure{le spectre}} \inserted{artinien}\handwritten{,} \deleted{\unsure{d'un corps},} \deleted{\unsure{d'un ann. artin|anneau artin}} \handwritten{ou} \deleted{\unsure{ou}} lorsque $S$ est un schéma local et $G$ est \handwritten{\textquotedblleft}réductif'' : dans ce cas, tout tore de $G$ est contenu dans un tore m\handwritten{a}ximal.

\uline{Définition} \handwritten{1}.5. Soit $G$ un groupe algébrique sur un corps $k$\handwritten{\unsure{.|,}} \deleted{\unsure{et soit}} On appelle rang réductif (resp. rang nilpotent, resp. rang unipotent ..) de $G$ le rang réductif (resp. ...) de $G_{\bar k}$, où $\bar k$ est une clôture algébrique de $k$.

On notera qu'en vertu de \handwritten{1}.2. et de la commutation de la formation de $\underline{\unsure{\mathbb C|C}\mathrm{entr}}_G(T)$ avec l'extension de la base, \deleted{\unsure{que}} les \deleted{rangs} notions de rang introduites dans \handwritten{1}.5. sont \uline{invariantes par extension du corps de base}; d'autre part, pour $k$ algébriquement clos, elles coïncident avec celles introduites au début du présent numéro.

\uline{Remarque} \handwritten{1}.6. Il n'est pas difficile de construire un schéma en groupes affine et lisse $G$ sur \deleted{\unsure{n}}le spectre $S$ d'un anneau de valuation discrète, dont la fibre générale soit isomorphe à $\underline G_m$, et la \deleted{\unsure{s}}fibre spéciale isomorphe à $\underline G_a$. Un tel $G$ ne contient aucun tore sauf le tore trivial \inserted{$T$} (réduit au sous-groupe unité), qui n'est évidemment pas tore maximal. De façon précise, dans la fibre spéciale $\unsure{\mathbb G|G}_{\unsure{o|0}}=\underline G_{a\ k}$, $T_o$ est bien tore maximal, mais dans la fibre générale $G_1=\underline G_{m\ K}$, $T_1$ n'est plus tore maximal (NB $k$ est le corps résiduel, $K$ le corps des fractions de l'anneau de valuation envisagé). On voit aussi sur cet exemple que le rang réductif de $G_s$ ($s\handwritten{\in}S$) n'est pas fonction continue de $s$. On a cependant les résultats suivants:

\uline{Théorème} \handwritten{1}.7. Soi\handwritten{e}nt $S$ un préschéma, $G$ un $S$-préschéma en groupes affine et lisse sur $S$. Pour tout $s\handwritten{\in}S$, considérons $\handwritten{\rho}_r(s)=\handwritten{\rho}_r(G_s)$ et $\handwritten{\rho}_n(s)=\handwritten{\rho}_n(G_s)$, les rangs réductif\deleted{\unsure{s}} et nilpotent de $G_s$ (\handwritten{1}.5.). Avec ces notations, on a ce qui suit:
% archive sheet number: 17
